\section{Seed Quality Assessment: Traditional Methods}

Traditional seed quality assessment primarily relies on manual visual inspection and laboratory testing. Visual inspection evaluates characteristics such as seed size, shape, colour, and texture, but the process is subjective and may vary between inspectors. Laboratory methods, including Tetrazolium (TZ) and germination tests, provide more accurate results but are time-consuming, destructive, and require controlled conditions.

Advanced techniques such as Near-Infrared (NIR) spectroscopy and X-ray imaging offer non-destructive seed analysis and can detect internal defects. However, their high cost and requirement for specialized equipment limit their adoption to well-equipped testing laboratories. As a result, farmers and local inspectors often lack access to rapid, affordable, and reliable seed verification methods.

\section{Deep Learning for Agricultural Image Detection}

Convolutional Neural Networks (CNNs) have been applied to a wide range of agricultural image detection tasks over the past decade, including leaf disease detection~\cite{mohanty2016plant}, weed detection~\cite{dyrmann2016}, and post-harvest quality assessment~\cite{zhang2020}. Their principal advantage over classical machine-learning approaches is the ability to learn spatial features directly from raw pixel data during training, eliminating the need for manually engineered feature extractors for each crop or defect type.

Applying CNNs to seed quality detection is a natural extension of this work. Paddy seeds often exhibit subtle variations in surface texture, colour, and visible defects between genuine and counterfeit batches. These differences are difficult to quantify using handcrafted features but can be effectively identified by CNNs when trained on sufficiently large datasets. Furthermore, the computational requirements of a trained CNN during inference are modest, making server-side web deployment feasible. This is particularly beneficial for farmers in rural areas who may have limited access to advanced hardware but can access the system through mobile internet connectivity.

\section{Transfer Learning and ResNet Architectures}

Training a CNN from random initialisation requires on the order of tens of thousands of labelled examples per class to achieve stable convergence. Annotated seed image datasets of that scale do not yet exist for most crop varieties, and certainly not for the specific task of detecting counterfeits. Transfer learning is the standard response: a model pre-trained on ImageNet (1.28 million labelled images across 1{,}000 categories) is fine-tuned on the smaller domain-specific dataset. The early convolutional layers, which encode low-level features such as edges and textures, are retained; only the later task-specific layers are updated. In practice, this approach consistently reaches accuracy close to what a from-scratch model would achieve on a dataset roughly ten times larger.

ResNet architectures are well-suited to this fine-tuning scenario. The residual skip connections introduced by He et al.~\cite{he2016deep} allow gradients to propagate through very deep networks without vanishing, making it feasible to fine-tune models with 50 or more layers on datasets that would cause shallower architectures to overfit. ResNet-50 has become a common baseline in agricultural image detection, balancing parameter count ($\approx$25 million) against accuracy on standard benchmarks, with pre-trained weights publicly available. For this project, it provides a well-validated starting point that reduces both training time and the risk of overfitting on the available paddy seed dataset.

\section{Farmer Feedback Systems}

A detection model trained on a fixed dataset will degrade over time if counterfeiting practices evolve and no new training data are collected. Field-level feedback from farmers is one practical mechanism for capturing these shifts. When a farmer submits a seed batch for image-based verification and subsequently reports germination failure or abnormal yield, that report constitutes a labelled example, a model prediction paired with a ground-truth outcome, that can be incorporated into periodic retraining cycles.

Feedback also serves a more immediate diagnostic function. Aggregating reports across users can flag a suspect lot or supplier before the model itself is updated: three independent failure reports from the same district within a fortnight is a meaningful signal, regardless of what the image detector predicts. This rule-based alerting complements the model rather than replacing it. Sustaining farmer participation remains a practical difficulty; deployments in low-resource contexts consistently find that single-tap rating interfaces produce far higher response rates than free-text fields or structured forms~\cite{munoz2018ict}.

\section{Web-Based Agricultural Platforms}

Deploying a trained model as a web service separates the computational workload from the end user's device. A farmer uploads a seed image through a browser; the server runs inference and returns a result. This architecture offers two concrete advantages over distributing a standalone mobile application: the model can be updated centrally without requiring client-side installations, and every prediction is logged in a structured format that supports monitoring and retraining.

For farmers in Manipur and comparable regions, browser-based access via a mid-range Android handset is more realistic than assuming a dedicated application installation. Keeping the frontend lightweight, avoiding heavy JavaScript bundles and large media assets reduces page-load times on 4G connections with variable signal quality. Containerised inference endpoints handle computationally intensive tasks without requiring a permanently running server, reducing operating costs during off-season periods when query volume is low.

\section{Review of Related Work}

Early automated seed detection systems employed Support Vector Machines (SVMs) and Random Forests using handcrafted features derived from shape measurements, colour histograms, and texture descriptors. Ni et al.~\cite{ni2009} identified rice seed varieties using morphological parameters extracted from binary silhouettes, achieving 94\% accuracy across five varieties under controlled lighting conditions. These methods perform adequately when variation within seed categories is low and imaging conditions are consistent, but they generalize poorly to photographs captured in real-world field environments.

Deep learning approaches have largely replaced handcrafted methods in recent studies. Liu et al.~\cite{liu2020} applied VGG16 transfer learning for maize seed defect detection, reporting 97.3\% accuracy on a 12-category dataset. Xu et al.~\cite{xu2021} used ResNet-50 for soybean quality assessment and found that fine-tuning from ImageNet weights outperformed training from scratch by approximately 8 percentage points when the training dataset contained fewer than 2{,}000 images per category. Both studies, however, focus on naturally occurring defects such as cracking, mould, and insect damage rather than intentional substitution.

A counterfeit seed is deliberately designed to resemble a genuine seed, making reliable identification significantly more challenging than detecting a damaged kernel. The visual differences between genuine and counterfeit seeds are often subtle, and the adaptive nature of counterfeiting means that counterfeit characteristics may evolve over time. To the best of our knowledge, no published study has directly evaluated CNN performance on verified counterfeit-versus-genuine paddy seed datasets, nor has any study integrated a structured farmer feedback mechanism into the seed verification process.

\section{Summary and Research Gaps}

The reviewed literature establishes that transfer-learned CNNs are effective for agricultural image detection, including seed quality assessment. Three gaps are directly relevant to this work. First, no published study has targeted intentional counterfeiting as distinct from natural quality defects; the two problems differ in that counterfeits may pass casual visual inspection and require the model to detect subtle discriminative cues. Second, existing models are evaluated on laboratory-quality images collected under controlled conditions; performance on photographs taken by farmers using mobile devices in variable lighting has not been reported. Third, no end-to-end system has been published that connects model inference to a structured feedback collection mechanism and uses that feedback iteratively for retraining.

This project addresses all three gaps. The proposed \textbf{AgriVerify} system is trained and evaluated on a counterfeit-versus-genuine paddy seed dataset. The web platform is designed to process images captured using mobile devices, making it accessible to farmers in real-world conditions. In addition, the farmer feedback module is integrated directly into the application, enabling users to report seed performance and field-level issues. These responses are stored in a structured format that can support future model improvement and system enhancement.
